% 05-safety.tex - Section 5: Safety Interface

\chapter{Safety Interface}
\label{sec:safety}

\section{Overview}

The SHL subsystem implements several safety mechanisms to protect equipment and personnel. These include temperature monitoring with automated responses, smoke and fire detection, water intrusion detection, and power monitoring.

\section{Temperature Monitoring}

\subsection{Temperature Thresholds}

The system monitors shelter temperature and takes action based on the configurable thresholds listed in Table~\ref{tab:temp-thresholds}.

\begin{table}[h]
    \centering
    \begin{tabular}{lll}
        \toprule
        \textbf{Condition} & \textbf{Default Threshold} & \textbf{Action} \\
        \midrule
        Normal & $<$ 80\textdegree F & No action \\
        Warning & 80--90\textdegree F & Set status to WARNING \\
        Critical & $\geq$ 90\textdegree F & Set status to ERROR, shutdown critical ports \\
        \bottomrule
    \end{tabular}
    \caption{Temperature Thresholds and Actions}
    \label{tab:temp-thresholds}
\end{table}

\begin{siteNote}
Threshold values are configurable per site via the \texttt{thermometers.warning\_temp} and \texttt{thermometers.critical\_temp} settings in \texttt{defaults.json}.
\end{siteNote}

\subsection{Critical Port Shutdown}

When the temperature exceeds the critical threshold, the system can automatically power off designated ``critical'' PDU ports. The list of critical ports is defined in the configuration file under \texttt{thermometers.critical\_list}.

Format: List of \texttt{[rack, port]} pairs.

Example:
\begin{verbatim}
"critical_list": [[1, 3], [1, 4], [2, 1]]
\end{verbatim}

\subsection{Temperature Status Messages}

The system \mib{INFO} field reflects the temperature status:
\begin{itemize}
    \item \texttt{TEMPERATURE! Shelter temperature at XX.X F, shutting down critical ports: ...}
    \item \texttt{TEMPERATURE! Shelter temperature at XX.X F} (warning level)
\end{itemize}

\section{Smoke and Fire Detection}

\begin{newInVersion}{F}
Available through a correctly configured EnviroMux.
\end{newInVersion}

When an EnviroMux system is configured with a smoke detector, the SHL subsystem monitors for smoke/fire conditions.

\subsection{Smoke Detection Response}

When smoke is detected:
\begin{enumerate}
    \item System status is set to \texttt{ERROR}
    \item \mib{INFO} is set to \texttt{SMOKE! Shelter smoke detector activated}
    \item Event is logged
\end{enumerate}

When smoke clears:
\begin{enumerate}
    \item System status returns to \texttt{NORMAL}
    \item \mib{INFO} indicates error condition cleared
\end{enumerate}

\begin{noteBox}
The smoke detector response is informational only. No automated equipment shutdown is triggered by smoke detection. Human intervention is required for fire response.
\end{noteBox}

\section{Water Intrusion Detection}

\begin{newInVersion}{F}
Available through a correctly configured EnviroMux.
\end{newInVersion}

When an EnviroMux system is configured with a water/leak sensor, the SHL subsystem monitors for water intrusion.

\subsection{Water Detection Response}

When water is detected:
\begin{enumerate}
    \item System status is set to \texttt{ERROR}
    \item \mib{INFO} is set to \texttt{WATER! Shelter water sensor activated}
    \item Event is logged
\end{enumerate}

When water clears:
\begin{enumerate}
    \item System status returns to \texttt{NORMAL}
    \item \mib{INFO} indicates error condition cleared
\end{enumerate}

\section{Door Monitoring}

\begin{newInVersion}{F}
Available through a correctly configured EnviroMux.
\end{newInVersion}

When an EnviroMux system is configured with a door sensor, the SHL subsystem monitors the shelter door status.

\subsection{Door Open Response}

To avoid false alarms during maintenance, the door warning is delayed:
\begin{itemize}
    \item Door open $<$ 4 hours: No warning
    \item Door open $\geq$ 4 hours: Set status to \texttt{WARNING}, \mib{INFO} indicates door is open
\end{itemize}

When the door is closed:
\begin{enumerate}
    \item Warning condition is cleared
    \item Door open timer is reset
\end{enumerate}

\section{Power Monitoring}

\begin{newInVersion}{F}
Available through the LWA line voltage monitor board.
\end{newInVersion}

The voltage monitor detects power quality issues on the AC mains.

\subsection{Power Flicker}

A power flicker is a brief voltage drop that recovers quickly. When detected:
\begin{enumerate}
    \item System status is set to \texttt{WARNING}
    \item \mib{INFO} includes \texttt{POWER-FLICKER! Flicker in the shelter power}
    \item Flicker status clears automatically after 5 minutes
\end{enumerate}

\subsection{Power Outage}

A power outage is a sustained loss of voltage. When detected:
\begin{enumerate}
    \item System status is set to \texttt{ERROR}
    \item \mib{INFO} is set to \texttt{POWER-OUTAGE! Shelter power outage}
    \item Outage state is saved to disk for persistence across restarts
    \item When power is restored, status returns to \texttt{NORMAL}
\end{enumerate}

\begin{noteBox}
The voltage monitor tracks both 120V and 240V lines independently. Either line experiencing an issue will trigger the corresponding warning or error.
\end{noteBox}

\section{Device Unreachable Monitoring}

The SHL subsystem monitors the reachability of all network devices (PDUs, thermometers, EnviroMux). When a device becomes unreachable:

\begin{enumerate}
    \item Device is added to the unreachable list
    \item If any devices are unreachable, status is set to \texttt{WARNING}
    \item \mib{INFO} shows count and list of unreachable devices
\end{enumerate}

Devices are removed from the unreachable list after successful communication is restored (with hysteresis to prevent flapping).

\section{Warning State Merging}

Multiple warning conditions can exist simultaneously. The system merges warning states using the format:

\begin{verbatim}
{MIB1}! {MIB2}! ... {description1}; {description2}; ...
\end{verbatim}

Example:
\begin{verbatim}
TEMPERATURE! DOOR! Shelter temperature at 82.0 F; Shelter door is open
\end{verbatim}

When a warning condition clears:
\begin{itemize}
    \item If other warnings remain: Update message to remove cleared condition
    \item If no warnings remain: Set status to \texttt{NORMAL}, set \mib{INFO} to ``Warning condition(s) cleared, system operating normally''
\end{itemize}
